\documentclass[10pt,twocolumn]{article}
\usepackage[utf8]{inputenc}
\usepackage[T1]{fontenc}
\usepackage[margin=1.9cm]{geometry}
\usepackage{times}
\usepackage{amsmath, amssymb}
\usepackage{booktabs}
\usepackage{graphicx}
\usepackage[numbers,sort&compress]{natbib}
\usepackage[hidelinks]{hyperref}

\newcommand{\model}{\textsc{BlockSparse}}
\DeclareMathOperator{\softmax}{softmax}

\title{Block-Sparse Attention for Long-Document Summarization}
\author{
  Razz\textsuperscript{1} \and Kenji Watanabe\textsuperscript{2} \and Aria Nakamura\textsuperscript{1} \\[4pt]
  \textsuperscript{1}Institute for Computational Linguistics \quad \textsuperscript{2}Department of Computer Science
}
\date{}

\begin{document}

\maketitle

\begin{abstract}
Transformer summarizers are limited by the quadratic cost of self-attention, which forces long documents to be truncated before they are read. We present \model{}, an attention pattern that combines a local window with a small set of learned block-level connections, and that runs in time linear in the document length. On three long-document benchmarks \model{} improves ROUGE-L by up to 1.8 points over truncation baselines while using 40\% less memory than full attention at 4{,}096 tokens. We release our code and pretrained checkpoints.
\end{abstract}

\section{Introduction}
\label{sec:introduction}

Abstractive summarization has improved rapidly since the introduction of the Transformer~\citep{vaswani2017attention} and of large pretrained sequence-to-sequence models~\citep{lewis2020bart}. Yet the documents people most want summarized, such as scientific articles, legal filings, and reports, are far longer than the 512 or 1{,}024 tokens these models accept. The usual remedy is to truncate the input, which discards most of the document.

The obstacle is self-attention, whose time and memory grow quadratically with the sequence length $n$. Sparse variants reduce this cost by restricting which pairs of tokens may attend to one another~\citep{beltagy2020longformer, zaheer2020bigbird}. Their global connections, however, are fixed in advance: a handful of special tokens, or a random subset.

We argue that for summarization the global pattern should follow the structure of the document. Our contributions are:
\begin{itemize}
  \item \model{}, an attention pattern in which each token attends to a local window and to a learned selection of block summaries;
  \item a differentiable top-$k$ selection that lets the pattern be trained end to end;
  \item experiments on arXiv, PubMed, and GovReport showing consistent gains over truncation and over fixed sparse patterns.
\end{itemize}

\section{Method}
\label{sec:method}

\subsection{Background}

Given queries $Q$, keys $K$, and values $V$ in $\mathbb{R}^{n \times d}$, scaled dot-product attention computes
\begin{equation}
  \operatorname{Attn}(Q, K, V) = \softmax\!\left( \frac{Q K^\top}{\sqrt{d}} \right) V ,
  \label{eq:attention}
\end{equation}
which needs $O(n^2)$ time and memory.

\subsection{Block summaries}

We partition the sequence into $B = \lceil n / b \rceil$ contiguous blocks of $b$ tokens and summarize block $j$ by the mean of its keys,
\begin{equation}
  \bar{k}_j = \frac{1}{b} \sum_{i \in \mathcal{B}_j} k_i .
\end{equation}
Each query $q_i$ scores every block by $s_{ij} = q_i^\top \bar{k}_j / \sqrt{d}$ and keeps the $r$ highest-scoring blocks. The token then attends to the tokens of its local window of width $w$ and to the summaries of the selected blocks.

\subsection{Complexity}

Scoring costs $O(nB)$, and attention costs $O(n(w + r))$. With $b = \sqrt{n}$ the total is $O(n\sqrt{n})$; with a fixed number of blocks it is $O(n)$. In all experiments we use $w = 256$, $b = 64$, and $r = 8$.

\subsection{Training}

Hard top-$r$ selection has zero gradient almost everywhere. During training we relax it with the perturbed-maximum estimator, adding Gumbel noise to the scores and annealing the temperature from $1.0$ to $0.1$ over the first 10{,}000 steps.

\section{Experiments}
\label{sec:experiments}

\subsection{Setup}

We initialize from BART-large~\citep{lewis2020bart}, replace encoder self-attention with \model{}, and extend the position embeddings to 4{,}096 tokens by copying. Models are fine-tuned for 50{,}000 steps with a batch size of 8 and a peak learning rate of $3 \times 10^{-5}$. We report ROUGE~\citep{lin2004rouge} on the test sets of arXiv, PubMed, and GovReport.

\subsection{Main results}

Table~\ref{tab:main} shows that reading more of the document helps on every dataset, and that learned block selection is better than a fixed global pattern at the same budget.

\begin{table}[t]
  \centering
  \caption{ROUGE-L on the test sets. Memory is the peak during training at batch size 1.}
  \label{tab:main}
  \small
  \begin{tabular}{lcccc}
    \toprule
    Model & arXiv & PubMed & GovRep. & Mem. \\
    \midrule
    BART (1{,}024) & 39.8 & 40.1 & 50.2 & 11\,GB \\
    Full attn. (4{,}096) & 41.5 & 41.9 & 53.0 & 38\,GB \\
    Longformer & 41.2 & 41.6 & 52.4 & 24\,GB \\
    BigBird & 41.3 & 41.8 & 52.7 & 26\,GB \\
    \model{} (ours) & \textbf{41.6} & \textbf{42.3} & \textbf{53.4} & 23\,GB \\
    \bottomrule
  \end{tabular}
\end{table}

\subsection{Ablations}

Removing the learned selection and attending to $r$ random blocks costs 0.7 ROUGE-L on arXiv. Halving the local window to $w = 128$ costs a further 0.4 points, which suggests that local and global context are complementary.

\section{Related Work}
\label{sec:related}

\paragraph{Sparse attention.} Longformer~\citep{beltagy2020longformer} combines a sliding window with a few global tokens, and BigBird~\citep{zaheer2020bigbird} adds random connections. Both fix the pattern before training; ours is learned per input.

\paragraph{Long-document summarization.} Hierarchical encoders summarize sections separately and then combine them, and extract-then-abstract pipelines first select salient sentences. \model{} keeps a single encoder and lets attention perform the selection.


\section{Conclusion}

We introduced \model{}, a sparse attention pattern for long documents whose global connections are learned rather than fixed. It matches or exceeds full attention on three summarization benchmarks at a fraction of the memory. In future work we plan to apply the same pattern to multi-document inputs.

\bibliographystyle{plainnat}
\bibliography{refs}

\end{document}
